\section{Conclusion and Recommendations}\label{sec:conclusion}
\subsection{Was the model good enough to design against?}
Yes. The gain and the hover voltage come from one regression that assumes neither, with an rms
residual of $0.022\ \mathrm{m\,s^{-1}}$ across eight runs spanning a sixfold range of step size,
and the fitted hover lands on the $2.50\ \mathrm{V}$ midpoint the brief says the simulator uses
without that value entering the fit. The time constant comes from $1800$ raw position samples and
is sharp enough that a $2$ percent error multiplies the residual by nine. Four routes to $A$ agree
inside $0.04$. The decay fit, which cannot hide an error in $A$ inside $T$, puts the time
constant $3.9$ percent below the step test, which is the one loose end in the identification and is
carried as such. That same fit predicts its own residual velocity from $A$ and $V_h$ to within
$0.008\ \mathrm{m\,s^{-1}}$, which is the strongest confirmation of the gain in the report.

The one disagreement is with the parameter-check tool of Appendix~\ref{sec:appendix}, which scored
the measured values at $54$ and $46$ percent against a target of $60$. It reports a number and no
reasoning, so it cannot say which of the two is wrong, and nothing in the data supports moving the
measurement toward it. The model was carried into the design unchanged.

\subsection{Did the controller meet the specifications?}
The verdicts are in Table~\ref{tab:spec_results}. The design behind them is $K_p = 0.005$ at a
$7\ \mathrm{m}$ setpoint with the bias trimmed to $2\ \mathrm{mV}$, predicted to overshoot
$17.3$ percent, hold a $6.0$ percent offset, and stay inside specification at every corner of both
robustness requirements.

\subsection{What limited the result}
Two limits are structural rather than experimental, and both would have applied to a perfect
build.

The first is that the two specifications pull against each other. Overshoot caps the gain through
$\zeta = 1/(2\sqrt{K_p k A T})$, and this plant has $kAT = 210$, so nothing above
$9.3\times 10^{-3}$ meets it, while the offset $\Delta V/(K_p k)$ is made large by that same small
gain. At a $1\ \mathrm{m}$ setpoint with a bias trimmed by eye the floor sits two orders of
magnitude above the ceiling and no proportional controller exists. What made the design possible
was reading the specification as written: the $8$ percent is a fraction of the setpoint and the
brief fixes none, so a larger setpoint relaxes the floor, with the sensor capping it from above.

The second is that proportional control cannot change the decay envelope. From
Equation~\eqref{eqn:wn_zeta}, $\zeta\omega_n = 1/(2T)$ contains no $K_p$: raising the gain raises
$\omega_n$ and makes the response ring without making it settle sooner. The settling time of
roughly $76\ \mathrm{s}$ belongs to the kitticopter, not the controller.

\subsection{Recommendations}
\begin{enumerate}
    \item \textbf{Add integral action.} An integrator drives the residual offset to zero and
    removes the dependence on the bias trim, currently the largest error source in the design.
    \item \textbf{Compensate rather than raise the gain.} With the envelope fixed at $1/(2T)$, a
    lead compensator is the only route to a faster response; more $K_p$ spends overshoot margin and
    buys no settling time.
    \item \textbf{Replace the bias divider with a reference.} The level shift is trimmed by hand
    off an unregulated rail, and Figure~\ref{fig:robustness_bias} shows the error is almost
    entirely proportional to its drift. A two-terminal reference removes both.
    \item \textbf{Match the difference amplifier resistors.} Four $10$ percent parts set the
    common-mode rejection of each difference stage, which no trimmer corrects.
\end{enumerate}
